\subsection{PSEUDOMETRICS}

DEFINITION 1. If X is a set, a pseudometric on X is any mapping f of X × X into the interval {[}0, +∞{]} of the extended real line R which satisfies the following conditions:

\((\mathrm{EC}_{\mathbf{I}})\) For all \(x \in \mathbf{X}, f(x, x) = 0\) .

\((\mathrm{EC}_{\mathrm{II}})\) For all \(x \in \mathbf{X}\) and all \(y \in \mathbf{X}, f(x, y) = f(y, x)\) (symmetry). \((\mathrm{EC}_{\mathrm{III}})\) For all \(x, y, z\) in \(\mathbf{X}\) ,

\[
f (x, y) \leqslant f (x, z) + f (z, y)
\]

(triangle inequality).

Examples. 1) On real number space \(\mathbf{R}^n\) , Euclidean distance (Chapter VI, § 2, no. 1) is a pseudometric.

\begin{enumerate}
\def\labelenumi{\arabic{enumi})}
\setcounter{enumi}{1}
\item
  If X is any set, the function f defined on \(X \times X\) by the conditions \(f(x, x) = 0\) for all \(x \in X\) , \(f(x, y) = +\infty\) if \(x \neq y\) is a pseudometric on X.
\item
  If \(\mathbf{X}\) is any set and if \(g\) is any finite real-valued function defined on \(\mathbf{X}\) , then the function \(f\) defined on \(\mathbf{X} \times \mathbf{X}\) by \(f(x, y) = |g(x) - g(y)|\) is a pseudometric on \(\mathbf{X}\) .
\end{enumerate}

* 4) Let X be the set of all continuous mappings of the interval {[}0, 1{]} of R into R. If for each pair of elements x, y of X we put

\[
f (x, y) = \int_ {0} ^ {1} | x (t) - y (t) | d t,
\]

then \(f\) is a pseudometric on \(\mathbf{X}\) . *

Remarks. 1) Example 2 above shows that a pseudometric can take the value +∞ for certain pairs of elements of X.

\begin{enumerate}
\def\labelenumi{\arabic{enumi})}
\setcounter{enumi}{1}
\tightlist
\item
  If \(f\) is a pseudometric on \(\mathbf{X}\) , we can in general have \(f(x, y) = 0\) without \(x = y\) , as Example 3 above shows (cf.~§ 2).
\end{enumerate}

From the triangle inequality it follows that if \(f(x, z)\) and \(f(y, z)\) are finite then so is \(f(x, y)\) ; moreover, in this case we have

\[
f (x, z) \leqslant f (y, z) + f (x, y) \quad \text { and } \quad f (y, z) \leqslant f (x, z) + f (x, y),
\]

so that

\[
\left| f (x, z) - f (y, z) \right| \leqslant f (x, y).\tag{1}
\]

If \(f\) is a pseudometric on \(\mathbf{X}\) , then so is \(\lambda f\) for any finite real number \(\lambda > 0\) . If \((f_t)_{t \in I}\) is any family of pseudometrics on \(\mathbf{X}\) , the sum \(\sum_{i \in I} f_i(x, y)\) is defined for all \((x, y) \in \mathbf{X} \times \mathbf{X}\) ; if \(f(x, y)\) denotes the value of this sum, then \(f\) is a pseudometric on \(\mathbf{X}\) . Again, the upper envelope \(g\) of the family \((f_t)\) (Chapter IV, § 5, no. 5) is a pseudometric on \(\mathbf{X}\) , for the relations \(f_i(x, y) \leqslant f_i(x, z) + f_i(z, y)\) imply

\[
\sup _ {t \in I} f _ {t} (x, y) \leqslant \sup _ {t \in I} \left(f _ {t} (x, z) + f _ {t} (z, y)\right) \leqslant \sup _ {t \in I} f _ {t} (x, z) + \sup _ {t \in I} f _ {t} (z, y)
\]

{[}Chapter IV, § 5, no. 7, formula (17){]}.

